\uuid{NOpf}
\titre{Coordonnées}
\auteur{bodin ; hanani ; mzari}
\organisation{exo7}
\licence{CC BY-NC-SA 4.0}
\datecreate{2020-01-20}
\niveau{L1}
\module{Géométrie}
\chapitre{Géométrie du plan}
\sousChapitre{Distance, norme, produit scalaire}
\difficulte{2}
\theme{coordonnées}
\exo7qcmid{lille-1.geometrie-plane.015}

\contenu{
\begin{question}
\qclassification{140.01,140.02}
\qtags{motcle=coordonnées}
On considère les points $A(1,1)$, $B(-1,1)$ et $C(1,-1)$. Quelles sont les bonnes réponses ?
\begin{answers}
    \bad{Les points $A$, $B$ et $C$ sont alignés.}
    \good{$ABC$ est un triangle rectangle en $A$.}
    \bad{$ABC$ est un triangle équilatéral.}
    \good{$ABC$ est un triangle isocèle en $A$.}
\end{answers}
\begin{explanations}
On a $\overrightarrow{AB}=(-2,0)$ et $\overrightarrow{AC}=(0,-2)$. Les points $A$, $B$ et $C$ ne sont pas alignés. De plus, $\|\overrightarrow{AB}\|=2=\|\overrightarrow{AC}\|$ donc $ABC$ est isocèle en $A$ et $\overrightarrow{AB}.\overrightarrow{AC}=0$, donc $ABC$ est rectangle en $A$.
\end{explanations}
\end{question}
}
